\documentclass[12pt]{article}
\usepackage{amsmath, amssymb}
\usepackage[margin=1in]{geometry}
\setlength{\parskip}{6pt}
\title{QUBO Formulation: Quantum-Circuit Qubit Allocation}
\date{}
\begin{document}
\maketitle

\subsection*{Quantum-Circuit Qubit Allocation}

\paragraph{Problem Statement.}
Given a set of logical qubits $L$, a set of physical qubits $P$, and a hardware coupling graph characterized by pairwise penalties $W_{p,q}$ for mapping logical interactions to physical pairs $(p,q)$, the objective is to determine a one-to-one assignment of logical qubits to physical qubits. The assignment must minimize the total penalty incurred by two-qubit gates, where $W_{p,q}$ typically encodes hardware distance or inverse gate fidelity. If $|L| > |P|$, the problem is infeasible.

\paragraph{Decision Variables.}
For each logical qubit $l \in L$ and physical qubit $p \in P$, introduce a binary variable $x_{l,p} \in \{0,1\}$, where $x_{l,p} = 1$ if logical qubit $l$ is mapped to physical qubit $p$, and $x_{l,p} = 0$ otherwise.

\paragraph{QUBO Derivation.}
\textbf{Step 1 -- Objective Function.} The original optimization problem minimizes the total allocation cost:
\begin{align}
\min_{x} \quad C(x) = \sum_{(i,j) \in E_{\text{log}}} \sum_{p \in P} \sum_{q \in P} W_{p,q} \, x_{i,p} \, x_{j,q},
\end{align}
where $E_{\text{log}}$ denotes the set of logical two-qubit gate edges. Diagonal terms $W_{p,p}$ are set to a sufficiently large constant to discourage self-mapping.

\textbf{Step 2 -- Constraints.} The assignment must satisfy two sets of constraints:
\begin{align}
\sum_{p \in P} x_{l,p} = 1 \quad \forall l \in L, \tag{C1}
\end{align}
\begin{align}
\sum_{l \in L} x_{l,p} \le 1 \quad \forall p \in P. \tag{C2}
\end{align}

\textbf{Step 3 -- Penalty Terms.} We convert the constraints into quadratic penalty terms using a positive weight $\lambda$.

For (C1), we use the standard squared penalty:
\begin{align}
P_1(x) = \lambda \sum_{l \in L} \left( \sum_{p \in P} x_{l,p} - 1 \right)^2.
\end{align}
Expanding the square and applying the idempotent property $x_{l,p}^2 = x_{l,p}$ yields:
\begin{align}
P_1(x) &= \lambda \sum_{l \in L} \left( \sum_{p \in P} x_{l,p}^2 + 2 \sum_{p < q} x_{l,p} x_{l,q} - 2 \sum_{p \in P} x_{l,p} + 1 \right) \nonumber \\
&= \lambda \sum_{l \in L} \left( 1 - \sum_{p \in P} x_{l,p} + 2 \sum_{p < q} x_{l,p} x_{l,q} \right).
\end{align}

For (C2), we use the inequality penalty form $\lambda \sum_p (\sum_l x_{l,p})(\sum_l x_{l,p} - 1)$, which naturally enforces the $\le 1$ condition without linear or constant terms:
\begin{align}
P_2(x) = \lambda \sum_{p \in P} \left( \sum_{l \in L} x_{l,p} \right) \left( \sum_{l \in L} x_{l,p} - 1 \right) = \lambda \sum_{p \in P} \sum_{l \neq l'} x_{l,p} x_{l',p}.
\end{align}

\textbf{Step 4 -- Combined Hamiltonian.} The full QUBO Hamiltonian $H(x)$ is the sum of the objective and penalty terms:
\begin{align}
H(x) = C(x) + P_1(x) + P_2(x).
\end{align}

\textbf{Step 5 -- Q Matrix and Offset.} Reading off the coefficients from $H(x)$, the QUBO matrix entries $Q_{(l,p),(l',p')}$ and constant offset $c$ are given by:
\begin{align}
Q_{(l,p),(l,p)} &= -\lambda, \\
Q_{(l,p),(l,q)} &= 2\lambda \quad \text{for } p \neq q, \\
Q_{(l,p),(l',p)} &= \lambda \quad \text{for } l \neq l', \\
Q_{(i,p),(j,q)} &= W_{p,q} \quad \text{for } (i,j) \in E_{\text{log}}, \\
c &= \lambda |L|.
\end{align}
All other entries are zero.

\paragraph{Penalty Weight Justification.}
Let $W_{\max} = \max_{p,q} |W_{p,q}|$ and let $d_{\max}$ be the maximum degree of any logical qubit in $E_{\text{log}}$. Violating a single constraint can at best reduce the objective cost by an amount bounded by $d_{\max} W_{\max}$. To guarantee that any infeasible assignment has strictly higher energy than the optimal feasible assignment, the penalty weight must satisfy:
\begin{align}
\lambda > d_{\max} \cdot W_{\max}.
\end{align}
In practice, a conservative choice such as $\lambda > 2 \max_{p,q} |W_{p,q}|$ or $\lambda = 100$ (for normalized $W_{p,q} \in [0,1]$) is sufficient to dominate any feasible objective improvement.

\paragraph{Correctness Argument.}
Minimizing $H(x)$ over $\{0,1\}^{|L| \times |P|}$ recovers the optimal allocation because $H(x) = C(x)$ for all feasible $x$, and $H(x) > C(x)$ for all infeasible $x$. Specifically, any violation of (C1) or (C2) introduces a penalty term strictly greater than $\lambda$. Since $\lambda$ is chosen to exceed the maximum possible reduction in $C(x)$ achievable by violating constraints, the global minimum of $H(x)$ cannot occur at an infeasible point. Consequently, the ground state of $H(x)$ corresponds exactly to a feasible mapping that minimizes the original allocation cost.

\paragraph{Illustrative Example.}
Consider a small instance with $|L| = 3$ logical qubits, $|P| = 3$ physical qubits, and logical edges $E_{\text{log}} = \{(0,1), (0,2), (1,2)\}$. Let all hardware penalties be uniform, $W_{p,q} = 1$ for $p \neq q$, and $W_{p,p} = 10$. Set $\lambda = 100$. The objective becomes:
\begin{align}
C(x) = \sum_{(i,j) \in \{(0,1),(0,2),(1,2)\}} \sum_{p=0}^{2} \sum_{q=0}^{2} W_{p,q} \, x_{i,p} x_{j,q}.
\end{align}
The penalty terms expand to:
\begin{align}
P_1(x) = 100 \sum_{l=0}^{2} \left( 1 - \sum_{p=0}^{2} x_{l,p} + 2 \sum_{0 \le p < q \le 2} x_{l,p} x_{l,q} \right),
\end{align}
\begin{align}
P_2(x) = 100 \sum_{p=0}^{2} \sum_{0 \le l < l' \le 2} x_{l,p} x_{l',p}.
\end{align}
The full Hamiltonian $H(x) = C(x) + P_1(x) + P_2(x)$ contains $3 \times 3 = 9$ binary variables. The constant offset is $c = 100 \times 3 = 300$. The optimal solution maps logical qubits $\{0,1,2\}$ to physical qubits $\{0,1,2\}$ in a permutation that aligns logical edges with the lowest physical penalties (e.g., $x_{0,0}=x_{1,1}=x_{2,2}=1$, yielding $C(x)=3$ and total energy $H(x)=303$). Any deviation from a valid permutation incurs a penalty of at least $100$, strictly increasing the energy.
\end{document}